\section{讨论与论断范围}
\label{sec:discussion}

UniGround 的首要科学贡献并非``通用控制器必须在每个宿主上胜过每种模型耦合方法''。而是：解码时接地控制可以外化到共享语义空间，并在异构 MLLM 家族上仍具操作层面的用处。因此讨论须通过两个耦合透镜解读经验结果：控制器是否真正可迁移，以及在该可迁移约束下其接地收益处于何种水平。

\subsection{何时通用主张成立}

仅当以下条件\textbf{同时}满足时，强通用主张才可成立。第一，同一 $\Psi_{\mathit{univ}}$ 检查点须在多个宿主模型上复用。第二，不得引入学习的按模型适配器。第三，通用路线须改善接地质量，或在不确定情形下足够频繁地弃权，以证明其计算开销合理。第四，当论文主张候选感知局部控制时，完整方法须实质优于更简单外部全局先验。缺少可迁移性、效用与比较必要性中的任一项，贡献表述就应收窄。

\subsection{不应混淆可迁移性与上限}

若 \texttt{TLRA\_internal\_calib} 在某一宿主上更强，这不应被当作通用表述的失败；它揭示预期的\textbf{可迁移性—上限权衡}：模型耦合标定因被允许利用内部几何，可能享有更高的宿主专用上限，而 UniGround 追求更可迁移的控制接口。论文应直接报告该权衡，而非将两路线强行纳入单一``优越性''叙事。

在此情形下，贡献收缩为更精确的陈述：UniGround 提供有原则的通用性框架、显式的可迁移性—上限分析，以及可复用的解码时控制器——即便模型耦合内部方法在单宿主上仍保留更高上界。

\subsection{何时须进一步收缩论断}

两类结果要求更强收缩。若完整通用控制器与 \texttt{External\_Global\_Prior} 或 \texttt{TLRA\_univ\_global\_only} 相当，则候选感知共振叙事尚不足以被认定为收益来源；论文应强调通过冻结外部证据实现的通用控制，但不宜强称 token 局部主张。若在附着时需要学习的按模型适配，则严格通用主张整体失效，须撤回。

这些收缩规则并非事后修辞防护；它们是论文科学叙事的一部分：方法旨在使通用命题可证伪，解释最终结果时也应保持这种可证伪性。
